\documentclass{report}
\usepackage{amsmath,amssymb}
\setlength{\parindent}{0mm}
\setlength{\parskip}{1em}
\begin{document}
\begin{center}
\rule{6in}{1pt} \
{\large Albert F. Lawrence \\
{\bf Iterative Methods in Large Field Electron Microscope Tomography }}

National Center for Microscopy and Imaging Research (NCMIR) \\ University of California \\ San Diego \\ 9500 Gilman Drive \\ La Jolla \\ CA 92093-0608 USA
\\
{\tt albert.rick.lawrence@gmail.com}\\
Xiaohua Wan\\
Sebastien Phan\\
Mark Ellisman\\
Fa Zhang\end{center}

Electron tomography is an emerging technology for the three dimensional
imaging of cellular ultrastructure. In combination with other techniques
this technology can provide three dimensional reconstructions of protein
assemblies, correlate 3D structures with functional investigations at the
light microscope level and provide structural information which extends
the findings of genomics and molecular biology. Because of rapid advances
in instrumentation and computer-based reconstruction, the routine imaging
of molecular structure in context appears to be likely within the next
few years.

Tomographic reconstruction from large format electron microscope images
requires special procedures to handle geometric distortions arising from
electron optics as opposed to light-ray optics. In particular, electrons
travel in curvilinear paths through the sample, defocus and other
aberrations can be more severe. We have developed a software package,
TxBR, which is based on a generalization of the inverse problem
associated with the ray transform. This software compensates for
instrument optics, sample degradation, and other deleterious effects.

TxBR utilizes a filtered backprojection algorithm for reconstruction, so
geometric compensation for curvilinear electron paths should also be
accompanied by a filtering algorithm which is accurate for the three
dimensional optics. This would require a pointwise treatment of the
filtering problem, which is, in effect, a six-dimensional kernel. At
present we employ an image warping technique which gives a usable
approximation to the correct kernel and reduces the problem to one
dimension. Further improvements would require alternative inversion
algorithms which bypass the filtering problem. Algorithms which employ
iterative techniques constitute an attractive class of alternatives. In
this case it is necessary to code corrections for the beam geometry into
the projections.

We present an adaptive simultaneous algebraic reconstruction technique
(ASART) in which a modified multilevel access scheme and an adaptive
relaxation parameter adjustment method are developed to improve the
quality of the reconstructed 3D structure. The reconstruction process is
facilitated by using a column-sum substitution approach. This modified
multilevel access scheme is adopted to arrange the order of projections
so as to minimize the correlations between consecutive views within a
limited angle range. In the adaptive relaxation parameter adjustment
method, not only the weight matrix (as in the existing methods) but the
gray levels of the pixels are employed to adjust the relaxation
parameters so that the quality of the reconstruction is improved while
the convergence process of the reconstruction is accelerated. In the
column-sum substitution approach, the computation to obtain the
reciprocal of the sum for the columns in each view is avoided so that the
needed computations for each iteration can be reduced.


\end{document}
